\documentclass{article}
\usepackage{graphicx}
\usepackage{amsmath}
\usepackage{cite}

\title{A Comprehensive Study on Multi-Task Learning for Diabetic Retinopathy and Diabetic Macular Edema Prediction from Retinal Fundus Images}
\author{RamaKrishna Sastry}
\date{\today}

\begin{document}

\maketitle

\begin{abstract}
    This research paper explores the utilization of multi-task learning techniques in the prediction of diabetic retinopathy (DR) and diabetic macular edema (DME) from retinal fundus images. We delve into the architectural details of the DR-ASPP-DRN model, its training methodology, and experimental outcomes, complemented by ablation studies and discussions on clinical significance.
\end{abstract}

\section{Introduction}
    \subsection{Problem Statement}
    Diabetic retinopathy (DR) and diabetic macular edema (DME) are leading causes of blindness. Early prediction is critical...\n    
    \subsection{Literature Review}
    Previous studies have explored various approaches for DR detection including...

\section{Architectural Details of DR-ASPP-DRN Model}
    The DR-ASPP-DRN model integrates...

\section{Methodology}
    The methodology employed includes...

\section{Datasets}
    Datasets used for training and validation include...

\section{Training Pipeline}
    The training pipeline is structured as follows...

\section{Experimental Results}
    Experimental results demonstrate...

\section{Ablation Studies}
    Ablation studies reveal the importance of...

\section{Challenges and Solutions}
    Key challenges faced during the research include...

\section{Clinical Significance}
    The clinical significance of these findings suggests...

\section{Conclusions and Future Work}
    In conclusion, the paper presents...

\bibliographystyle{plain}
\bibliography{references}

\end{document}